% ═══════════════════════════════════════════════════════════════
% SPANISCH LERNBLATT: de + artículo (del) + Así se dice
% ═══════════════════════════════════════════════════════════════
% Thema: Besitz ausdrücken + Klassenraumredemittel
% Klasse 7 | Unidad 2 | Mi instituto | DS 03
% ═══════════════════════════════════════════════════════════════

\documentclass[10pt, a4paper]{article}

\usepackage[utf8]{inputenc}
\usepackage[T1]{fontenc}
\usepackage[ngerman]{babel}

% --- SCHRIFTEN ---
\usepackage{palatino}
\usepackage{helvet}
\newcommand{\sanstext}[1]{{\fontfamily{phv}\selectfont #1}}

\usepackage{microtype}
\usepackage[margin=1.4cm, top=1.6cm, bottom=1.3cm]{geometry}
\usepackage{enumitem}
\usepackage{tabularx}
\usepackage{booktabs}
\usepackage{array}
\usepackage{multicol}
\usepackage{tikz}
\usetikzlibrary{calc}

\usepackage{fancyhdr}
\usepackage{lastpage}
\usepackage{needspace}

\usepackage{xcolor}
\usepackage{amssymb}
\usepackage{tcolorbox}
\tcbuselibrary{skins}

% ═══════════════════════════════════════════════════════════════
% FARBEN
% ═══════════════════════════════════════════════════════════════

\definecolor{kopfzeile}{HTML}{1A1A1A}
\definecolor{akzent}{HTML}{C41E3A}          % Spanisch-Karminrot
\definecolor{hellgrau}{HTML}{F2F2F2}
\definecolor{mittelgrau}{HTML}{777777}
\definecolor{dunkelgrau}{HTML}{444444}

% ═══════════════════════════════════════════════════════════════
% VARIABLEN
% ═══════════════════════════════════════════════════════════════

\newcommand{\thema}{de + artículo (del)}
\newcommand{\klasse}{7}

% ═══════════════════════════════════════════════════════════════
% KOPF- UND FUSSZEILE
% ═══════════════════════════════════════════════════════════════

\pagestyle{fancy}
\fancyhf{}
\renewcommand{\headrulewidth}{0pt}
\renewcommand{\footrulewidth}{0pt}
\cfoot{\sanstext{\footnotesize\textcolor{mittelgrau}{\thepage\,/\,\pageref{LastPage}}}}

% ═══════════════════════════════════════════════════════════════
% BOXEN
% ═══════════════════════════════════════════════════════════════

% Grammatik-Box (dunkelgrauer Strich oben)
\newtcolorbox{grammatikbox}[1][]{
    colback=hellgrau,
    colframe=hellgrau,
    boxrule=0pt,
    arc=0pt,
    left=8pt, right=8pt, top=8pt, bottom=6pt,
    borderline north={1.5pt}{0pt}{dunkelgrau},
    title={\sanstext{\small\bfseries #1}},
    coltitle=dunkelgrau,
    attach boxed title to top left={yshift=-2pt, xshift=8pt},
    boxed title style={colback=hellgrau, colframe=hellgrau, boxrule=0pt, arc=0pt}
}

% Merk-Box (roter Strich oben)
\newtcolorbox{merkbox}[1][]{
    colback=hellgrau,
    colframe=hellgrau,
    boxrule=0pt,
    arc=0pt,
    left=8pt, right=8pt, top=8pt, bottom=6pt,
    borderline north={1.5pt}{0pt}{akzent},
    title={\sanstext{\small\bfseries\textcolor{akzent}{#1}}},
    coltitle=akzent,
    attach boxed title to top left={yshift=-2pt, xshift=8pt},
    boxed title style={colback=hellgrau, colframe=hellgrau, boxrule=0pt, arc=0pt}
}

% Vokabel-Box (dünner Rahmen)
\newtcolorbox{vokabelbox}[1][]{
    colback=white,
    colframe=mittelgrau,
    boxrule=0.5pt,
    arc=0pt,
    left=8pt, right=8pt, top=8pt, bottom=6pt,
    title={\sanstext{\small\bfseries #1}},
    coltitle=dunkelgrau,
    attach boxed title to top left={yshift=-2pt, xshift=8pt},
    boxed title style={colback=white, colframe=white, boxrule=0pt, arc=0pt}
}

% AB-Verweis
\newcommand{\abmark}[1]{%
    \tikz[baseline=(X.base)]\node[fill=akzent, text=white,
    font=\sanstext\tiny\bfseries, inner sep=2pt, rounded corners=1pt](X){→ AB #1};%
}

% ═══════════════════════════════════════════════════════════════
% BEFEHLE
% ═══════════════════════════════════════════════════════════════

% Spanisch: Palatino fett
\newcommand{\es}[1]{\textbf{#1}}

% Deutsch: klein, kursiv, grau
\newcommand{\de}[1]{{\small\textit{\textcolor{mittelgrau}{#1}}}}

% Grammatikbegriff: Kapitälchen
\newcommand{\gram}[1]{\textsc{\small #1}}

\setlength{\parindent}{0pt}
\setlength{\parskip}{4pt}
\setlength{\columnsep}{20pt}

% ═══════════════════════════════════════════════════════════════
% DOKUMENT
% ═══════════════════════════════════════════════════════════════

\begin{document}

% ═══════════════════════════════════════════════════════════════
% HEADER
% ═══════════════════════════════════════════════════════════════

\noindent
\begin{tikzpicture}[remember picture, overlay]
    \fill[kopfzeile] (current page.north west) rectangle +(\paperwidth, -1cm);
    \node[anchor=west, text=white, font=\sanstext\large\bfseries]
        at ([xshift=1.4cm, yshift=-0.5cm]current page.north west) {\thema};
    \node[anchor=east, text=white, font=\sanstext\small]
        at ([xshift=-1.4cm, yshift=-0.5cm]current page.north east)
        {Spanisch\,·\,Klasse \klasse\,·\,Lernblatt};
\end{tikzpicture}

\vspace{0.6cm}

% ═══════════════════════════════════════════════════════════════
% INHALT
% ═══════════════════════════════════════════════════════════════

\begin{multicols}{2}

% ═══════════════════════════════════════════════════════════════
% ABSCHNITT 1: de + artículo
% ═══════════════════════════════════════════════════════════════

\begin{grammatikbox}[de + artículo]

\vspace{2pt}

Mit \es{de} \de{von} drückst du Besitz oder Zugehörigkeit aus:

\vspace{4pt}

{\small
\begin{tabular}{@{}lcl@{\hspace{8pt}}l@{}}
de + el & $\rightarrow$ & \es{del} & \de{Kontraktion!} \\
de + la & $\rightarrow$ & de la & \\
de + los & $\rightarrow$ & de los & \\
de + las & $\rightarrow$ & de las & \\
\end{tabular}}

\vspace{6pt}

\textbf{Beispiele:}

\vspace{2pt}

{\small
\begin{itemize}[leftmargin=1em, itemsep=2pt]
    \item el libro \es{del} profesor \de{das Buch des Lehrers}
    \item la silla \es{de la} profesora \de{der Stuhl der Lehrerin}
    \item los cuadernos \es{de los} alumnos \de{die Hefte der Schüler}
    \item la mochila \es{de las} alumnas \de{der Rucksack der Schülerinnen}
\end{itemize}}

\end{grammatikbox}

\vspace{4pt}

\abmark{1-2}

\vspace{8pt}

% ═══════════════════════════════════════════════════════════════
% ABSCHNITT 2: Repaso hay vs estar
% ═══════════════════════════════════════════════════════════════

\begin{merkbox}[Repaso: hay vs estar]

\vspace{2pt}

{\small
\textbf{HAY} = \textit{Es gibt} (Existenz, unbestimmt)

\begin{itemize}[leftmargin=1em, itemsep=1pt]
    \item + un/una: \es{Hay un libro.}
    \item + Zahl: \es{Hay tres sillas.}
    \item + mucho/poco: \es{Hay muchos cuadernos.}
\end{itemize}

\vspace{4pt}

\textbf{ESTAR} = \textit{Wo ist?} (Ort, bestimmt)

\begin{itemize}[leftmargin=1em, itemsep=1pt]
    \item + el/la: \es{El libro está en la mesa.}
    \item + mi/tu: \es{Mi mochila está aquí.}
\end{itemize}}

\vspace{4pt}

{\footnotesize\textbf{Tipp:} Frag dich: \textit{Gibt es...?} → hay | \textit{Wo ist...?} → estar}

\end{merkbox}

\vspace{4pt}

\abmark{3}

\columnbreak

% ═══════════════════════════════════════════════════════════════
% ABSCHNITT 3: Así se dice
% ═══════════════════════════════════════════════════════════════

\begin{vokabelbox}[Así se dice – Im Unterricht]

\vspace{2pt}

{\small
\renewcommand{\arraystretch}{1.3}
\begin{tabular}{@{}p{4.2cm}l@{}}
\es{¿Cómo?} & Wie bitte? \\
\es{No comprendo.} & Ich verstehe nicht. \\
\es{Tengo una pregunta.} & Ich habe eine Frage. \\
\es{¿Cómo se dice ... en español?} & Wie sagt man ... auf Spanisch? \\
\es{Más despacio, por favor.} & Langsamer, bitte. \\
\es{¿Puedes repetir?} & Kannst du wiederholen? \\
\es{¿Qué significa ...?} & Was bedeutet ...? \\
\end{tabular}}

\end{vokabelbox}

\vspace{4pt}

\abmark{4}

\vspace{8pt}

% ═══════════════════════════════════════════════════════════════
% ABSCHNITT 4: Übersicht Artikel
% ═══════════════════════════════════════════════════════════════

\begin{grammatikbox}[Übersicht: Bestimmte Artikel]

\vspace{2pt}

{\small
\renewcommand{\arraystretch}{1.2}
\begin{tabular}{@{}l|cc@{}}
 & \textbf{Singular} & \textbf{Plural} \\
\hline
\textbf{maskulin} & el & los \\
\textbf{feminin} & la & las \\
\end{tabular}}

\vspace{6pt}

\textbf{Mit de:}

{\small
\renewcommand{\arraystretch}{1.2}
\begin{tabular}{@{}l|cc@{}}
 & \textbf{Singular} & \textbf{Plural} \\
\hline
\textbf{maskulin} & \es{del} (de+el) & de los \\
\textbf{feminin} & de la & de las \\
\end{tabular}}

\vspace{4pt}

{\footnotesize\textbf{Merke:} Nur \es{de + el = del} wird zusammengezogen!}

\end{grammatikbox}

\vspace{8pt}

% ═══════════════════════════════════════════════════════════════
% ABSCHNITT 5: Spiel
% ═══════════════════════════════════════════════════════════════

\begin{merkbox}[Spiel: ¿De quién es?]

\vspace{2pt}

{\small
\textbf{So geht's:}
\begin{enumerate}[leftmargin=1.5em, itemsep=1pt]
    \item Jeder legt 1-2 Gegenstände auf den Tisch.
    \item Der Lehrer zeigt einen Gegenstand: \es{``¿De quién es este bolígrafo?''}
    \item Ihr ratet: \es{``Es el bolígrafo de María.''} oder \es{``Es del profesor.''}
    \item Der Besitzer sagt: \es{``¡Sí, es mi bolígrafo!''} oder \es{``No, no es mi bolígrafo.''}
\end{enumerate}

\vspace{4pt}

\textbf{Redemittel:}
\begin{itemize}[leftmargin=1em, itemsep=1pt]
    \item Es el/la ... \es{de} [Name].
    \item Es el/la ... \es{del} profesor / \es{de la} profesora.
    \item Es el/la ... \es{de los} alumnos / \es{de las} alumnas.
\end{itemize}}

\end{merkbox}

\end{multicols}

% ═══════════════════════════════════════════════════════════════
% FOOTER
% ═══════════════════════════════════════════════════════════════

\vfill

\noindent\rule{\textwidth}{0.5pt}

\vspace{2pt}

{\footnotesize\textcolor{mittelgrau}{Qué pasa 1, S. 42-43 | Unidad 2: Mi instituto}}

\end{document}
